\DocumentMetadata{
  pdfversion=2.0,
  pdfstandard=ua-2,
  lang=en-US,
  tagging=on,
  tagging-setup={math/setup=mathml-SE,tabsorder=structure}
}
\documentclass[11pt,letterpaper]{article}
\usepackage[margin=0.68in,includefoot]{geometry}
\usepackage{newcomputermodern}
\usepackage{amsmath,amscd,mathtools}
\usepackage{tikz,adjustbox}
\providecommand{\square}{\mdlgwhtsquare}
\usepackage[hidelinks]{hyperref}
\hypersetup{
  pdftitle={Analysis PhD Exam, September 2010},
  pdfauthor={University of Florida Department of Mathematics},
  pdfsubject={Graduate mathematics examination},
  pdfdisplaydoctitle=true,
  bookmarksopen=true
}
\setcounter{secnumdepth}{0}
\setlength{\parindent}{0pt}
\setlength{\parskip}{0.28em}
\setlength{\emergencystretch}{3em}
\setlength{\leftmargini}{2.15em}
\setlength{\leftmarginii}{2.65em}
\setlength{\leftmarginiii}{3em}
\renewcommand{\labelenumi}{\textbf{\theenumi.}}
\pagestyle{empty}
\raggedbottom
\allowdisplaybreaks
\newenvironment{examproblems}[1]{%
  \begin{enumerate}%
  \setcounter{enumi}{#1}%
  \setlength{\topsep}{0.35em}%
  \setlength{\partopsep}{0pt}%
  \setlength{\itemsep}{0.48em}%
  \setlength{\parsep}{0.18em}%
}{\end{enumerate}}
\newenvironment{examsubparts}{%
  \begin{enumerate}%
  \setlength{\topsep}{0.18em}%
  \setlength{\partopsep}{0pt}%
  \setlength{\itemsep}{0.16em}%
  \setlength{\parsep}{0.1em}%
}{\end{enumerate}}
\newenvironment{examinstructions}{%
  \par\begingroup\itshape
}{%
  \par\endgroup\vspace{0.18em}
}
\makeatletter
\renewcommand\section{\@startsection{section}{1}{0pt}%
  {-1.25ex plus -.25ex minus -.1ex}%
  {0.55ex plus .1ex}%
  {\normalfont\large\bfseries}}
\renewcommand{\@maketitle}{%
  \newpage\null\vskip -1em%
  {\footnotesize\bfseries
    \href{https://www.ufl.edu/}{University of Florida}, \href{https://math.ufl.edu/}{Department of Mathematics}\par}%
  \vskip -0.5em%
  \tagstructbegin{tag=Title}%
    {\LARGE\bfseries\href{https://gma.math.ufl.edu/exams/analysis-phd/}{Analysis PhD Exam}, September 2010\par}%
  \tagstructend%
  \vskip 0.25em%
  \tagmcbegin{artifact}%
    \hrule height 1pt%
  \tagmcend%
  \vskip 1em%
}
\makeatother
\title{Analysis PhD Exam, September 2010}
\author{}
\date{}
\begin{document}
\maketitle
\thispagestyle{empty}
\begin{examinstructions}
Write solutions in a neat and logical fashion, giving complete reasons for all steps.

\par
Attempt EIGHT problems.
\end{examinstructions}
\begin{examproblems}{0}
\item
Let \(f\in L^1[0,1]\). Suppose that 
\[
\int_0^1 x^n f(x)\,dx=0 \quad\text{for all } n=0,1,2,\ldots
\]
 Prove that \(f=0\) almost everywhere. (Hint: reduce the problem to showing that if \(\int_E f(x)\,dx=0\) for every Borel~\(E\), then \(f=0\) a.e.)
\item
Let \((X,\mathcal M,\mu)\) be a~\(\sigma\)-finite measure space, and let \(g:X\to\mathbb C\) be a measurable function. Prove that the operator \(M_g:f\mapsto gf\) is bounded on \(L^p(\mu)\) (\(1\leq p\leq\infty\)) if and only if \(g\in L^\infty(\mu)\), in which case \(\lVert M_g\rVert=\lVert g\rVert_\infty\).
\item
This problem has two parts.
\begin{examsubparts}
\item[\textnormal{a)}]
Define monotone class, and state the Monotone Class Theorem.
\item[\textnormal{b)}]
State the Fubini–Tonelli theorem.
\end{examsubparts}
\item
This problem has two parts.
\begin{examsubparts}
\item[\textnormal{a)}]
Let~\(\nu\) be a signed measure on \((X,\mathcal M)\). Define what it means for a set to be positive, negative or null for~\(\nu\), and state the Hahn decomposition theorem.
\item[\textnormal{b)}]
Prove that if~\(\nu\) is a signed measure, then there exist unique positive measures \(\nu^+,\nu^-\) so that \(\nu=\nu^+-\nu^-\) and \(\nu^+\perp\nu^-\).
\end{examsubparts}
\item
Let \((k_n)_{n=1}^\infty\) be a sequence of measurable, \(1\)-periodic functions on \(\mathbb R\) with the following properties:

\par
Prove that if \(f\in L^1[-\tfrac12,\tfrac12]\), then \(f*k_n\to f\) in \(L^1[-\tfrac12,\tfrac12]\).

\par
(Here \(*\) denotes convolution on \([-\tfrac12,\tfrac12]\): \((f*g)(x)=\int_{-1/2}^{1/2}f(t)g(x-t)\,dt\).)
\begin{examsubparts}
\item[\textnormal{a)}]
\(k_n\geq 0\) a.e., for all~\(n\).
\item[\textnormal{b)}]
\(\displaystyle \int_{-1/2}^{1/2}k_n(t)\,dt=1\) for all~\(n\).
\item[\textnormal{c)}]
For each \(\delta>0\), \(k_n\to 0\) uniformly on the set \(\delta\leq |t|\leq\tfrac12\).
\end{examsubparts}
\end{examproblems}
\begin{examinstructions}
A second page was presumably part of the original exam, but it was not present in the exam archive.
\end{examinstructions}
\end{document}
